    \section{Assignment Overview \\
    \small{\textit{-- Joris Wilson, Eva Kemp, and Nica Saoi}}}
    \index{Section!Assignment Overview}
    \label{Section::Assignment Overview}

    This assignment extends the FTDI oscilloscope project with a Qt-based user interface while preserving the existing MVC architecture, the scale-and-shift processing pipeline, and the multithreaded FTDI transfer path. The current implementation provides two selectable Qt user interfaces, keeps one abstract View boundary, and routes live acquisition through the existing circular-buffer pipeline instead of moving hardware logic into the UI layer.

    \section{Requirements \\
    \small{\textit{-- Joris Wilson, Eva Kemp, and Nica Saoi}}}
    \index{Section!Requirements}
    \label{Section::Requirements}

    \begin{enumerate}
        \item The system shall provide two selectable user interfaces.
        \item The system shall support pipeline processing for scale and shift operations.
        \item The system shall read from and write to the FTDI device using multithreaded circular buffers.
        \item Extra credit: The system shall write data to one FTDI device and read data from a second FTDI device using two circular buffers, one for writing and one for reading.
        \item The implementation shall use one class per file.
        \item The system architecture shall be documented with at least a class diagram for this assignment.
        \item The system shall use the FTDI libraries to access the FTDI chips.
        \item For this assignment, namespaces shall follow the format \texttt{io<ClassName>}.
    \end{enumerate}

    \section{Architecture Overview \\
    \small{\textit{-- Joris Wilson, Eva Kemp, and Nica Saoi}}}
    \index{Section!Architecture Overview}
    \label{Section::Architecture Overview}

    \subsection{Architectural Style}
    The primary user-interface path follows the Model--View--Controller architectural style. The View layer is represented by one abstract Qt-facing boundary, the Controller layer coordinates user intent and live-session behavior, and the Model layer owns waveform state, control state, viewport state, and the active signal-processing pipeline.

    \subsection{MVC Structure}
    The controller layer is represented by \texttt{ioOscilloscopeController}. It coordinates UI actions, model updates, render-state generation, and live acquisition startup and shutdown.

    The model layer is represented by \texttt{ioOscilloscopeModel}. It stores raw samples, processed samples, control state, viewport state, and the current filter pipeline.

    The view layer is represented by the abstract class \texttt{ioOscilloscopeView}. The two selectable Qt user interfaces are the concrete subclasses \texttt{ioCompactOscilloscopeView} and \texttt{ioDetailedOscilloscopeView}.

    The class \texttt{ioQtScopeWindow} is the top-level Qt host window. It switches between concrete view instances and forwards user selections to the controller, but it is not a second View abstraction.

    \subsection{Two Selectable User Interfaces}
    The current implementation provides two selectable Qt user interfaces. The first is \texttt{ioCompactOscilloscopeView}, which emphasizes a compact, controls-first layout and a concise waveform status panel.

    The second is \texttt{ioDetailedOscilloscopeView}, which emphasizes a larger waveform region and a telemetry-oriented detail panel. Both views inherit from \texttt{ioOscilloscopeView}, so the controller renders through one shared View contract.

    \subsection{Pipeline Processing Structure}
    Scale and shift operations remain in a dedicated pipe-and-filter path rather than being embedded in the user interface. The model owns an \texttt{ioFilterPipeline}, and the controller updates the model when scale or offset changes.

    The pipeline is built from \texttt{ioSignalFilter}, \texttt{ioSampleMappingFilterBase}, \texttt{ioScaleFilter}, and \texttt{ioOffsetFilter}. The filters are applied in deterministic order so the processed waveform is reproducible and easy to reason about.

    \subsection{FTDI Communication Structure}
    FTDI access is isolated inside \texttt{ioLibrary}. The low-level D2XX session boundary is owned by \texttt{ioFtdiSession}, while the stream-level abstractions \texttt{ioFtdiByteStream} and \texttt{ioFtdiOutputByteStream} expose readable and writable FTDI endpoints to the rest of the system.

    The adapter \texttt{ioFtdiDevice} is defined in \texttt{ftd2xx\_wrapper.py}. It preserves a simple device-oriented API while reusing the canonical session layer.

    \subsection{Multithreaded Circular Buffer Structure}
    The required multithreaded transfer path remains outside the UI layer. The class \texttt{ioPipelineController} builds and coordinates a producer-consumer pipeline composed of \texttt{ioUsbReadController}, \texttt{ioDataBuffer}, and \texttt{ioUsbWriteController}.

    The class \texttt{ioDataBuffer} is the required bounded thread-safe circular buffer. It supports blocking \texttt{push()} and \texttt{pop()} operations, bounded capacity, synchronized full and empty waiting, and controlled close behavior for safe shutdown.

    \subsection{Live Session Structure}
    The oscilloscope's live acquisition path is owned by \texttt{ioLiveOscilloscopeSession}. This class creates the transfer pipeline, captures rolling waveform history, and keeps FTDI or file-backed sampling outside the View layer.

    The live-session helpers \texttt{ioTappedReadableByteStream}, \texttt{ioLiveSampleHistory}, \texttt{ioLiveFileTailByteStream}, \texttt{ioGeneratedWaveformByteStream}, and \texttt{ioNullWritableByteStream} support live acquisition and tee-output behavior without collapsing MVC boundaries.

    \section{Class List and Responsibilities \\
    \small{\textit{-- Joris Wilson, Eva Kemp, and Nica Saoi}}}
    \index{Section!Class List and Responsibilities}
    \label{Section::Class List and Responsibilities}

    \subsection{Bootstrap and CLI Classes}
    \begin{itemize}
        \item \texttt{ioScopeShell} \\
        Responsibility: Runs the oscilloscope shell and persists runtime settings.

        \item \texttt{ioPipelineCli} \\
        Responsibility: Builds transfer configuration and runs the multithreaded FTDI or file pipeline from the command line.

        \item \texttt{ioLegacyFtdiCli} \\
        Responsibility: Preserves older direct FTDI demo commands for low-level chip interaction.

        \item \texttt{ioFileComparator} \\
        Responsibility: Compares two binary files byte for byte for deterministic transfer verification.
    \end{itemize}

    \subsection{Qt MVC Classes}
    \begin{itemize}
        \item \texttt{ioOscilloscopeController} \\
        Responsibility: Coordinates UI actions, model updates, render-state fan-out, and live-session control.

        \item \texttt{ioOscilloscopeModel} \\
        Responsibility: Stores waveform data, viewport state, control state, and the active processing pipeline.

        \item \texttt{ioOscilloscopeView} \\
        Responsibility: Defines the single abstract View contract used by both selectable Qt user interfaces.

        \item \texttt{ioCompactOscilloscopeView} \\
        Responsibility: Implements the compact controls-first Qt user interface.

        \item \texttt{ioDetailedOscilloscopeView} \\
        Responsibility: Implements the detailed waveform-first Qt user interface.

        \item \texttt{ioQtScopeWindow} \\
        Responsibility: Hosts the selectable Qt views and forwards top-level UI selections to the controller.

        \item \texttt{ioOscilloscopeWindow} \\
        Responsibility: Defines the generic window-host protocol used outside the primary Qt path.

        \item \texttt{ioTkOscilloscopeWindow} \\
        Responsibility: Preserves the older Tk-based oscilloscope window as a legacy compatibility path.

        \item \texttt{ioControlState} \\
        Responsibility: Stores sampling, source-selection, FTDI, scale, offset, tee-output, and active-view settings.

        \item \texttt{ioViewportState} \\
        Responsibility: Stores the visible waveform window position and size.

        \item \texttt{ioRenderState} \\
        Responsibility: Carries immutable render data from the controller into each concrete view.

        \item \texttt{ioScopeSettingsStore} \\
        Responsibility: Loads and saves the oscilloscope settings JSON payload.
    \end{itemize}

    \subsection{Live Session Support Classes}
    \begin{itemize}
        \item \texttt{ioLiveOscilloscopeSession} \\
        Responsibility: Owns the live acquisition session and its tee-output pipeline.

        \item \texttt{ioLiveSampleHistory} \\
        Responsibility: Stores rolling normalized samples for live waveform rendering.

        \item \texttt{ioTappedReadableByteStream} \\
        Responsibility: Mirrors live input bytes into sample history while preserving readable-stream behavior.

        \item \texttt{ioLiveFileTailByteStream} \\
        Responsibility: Tails a file as a continuously readable byte stream.

        \item \texttt{ioGeneratedWaveformByteStream} \\
        Responsibility: Exposes generated demo waveforms as a readable byte stream.

        \item \texttt{ioNullWritableByteStream} \\
        Responsibility: Discards tee output while preserving pipeline flow.
    \end{itemize}

    \subsection{Signal Pipeline Classes}
    \begin{itemize}
        \item \texttt{ioSignalFilter} \\
        Responsibility: Defines the interface for one waveform-processing stage.

        \item \texttt{ioSampleMappingFilterBase} \\
        Responsibility: Shares one-sample-at-a-time transform behavior for simple filters.

        \item \texttt{ioScaleFilter} \\
        Responsibility: Applies the vertical scale transformation.

        \item \texttt{ioOffsetFilter} \\
        Responsibility: Applies the vertical shift transformation.

        \item \texttt{ioFilterPipeline} \\
        Responsibility: Applies the configured waveform filters in deterministic order.
    \end{itemize}

    \subsection{Signal Source and Theme Classes}
    \begin{itemize}
        \item \texttt{ioSignalSource} \\
        Responsibility: Defines the strategy contract for waveform acquisition.

        \item \texttt{ioGeneratedWaveformSource} \\
        Responsibility: Produces deterministic generated demo waveforms.

        \item \texttt{ioFileWaveformSource} \\
        Responsibility: Produces normalized samples from file input.

        \item \texttt{ioFtdiWaveformSource} \\
        Responsibility: Produces normalized samples from FTDI device input.

        \item \texttt{ioWaveformGenerator} \\
        Responsibility: Selects the active waveform source from the current control state.

        \item \texttt{ioViewTheme} \\
        Responsibility: Defines the theme contract for a view.

        \item \texttt{ioStaticThemeBase} \\
        Responsibility: Shares static theme metadata behavior.

        \item \texttt{ioPortraitTheme} \\
        Responsibility: Supplies the portrait-oriented palette used by the compact view.

        \item \texttt{ioLandscapeTheme} \\
        Responsibility: Supplies the landscape-oriented palette used by the detailed view.
    \end{itemize}

    \subsection{FTDI Stream and Session Classes}
    \begin{itemize}
        \item \texttt{ioFtdiDevice} \\
        Responsibility: Provides the FTDI device adapter over the session layer.

        \item \texttt{ioLibraryError} \\
        Responsibility: Defines the base runtime error for \texttt{ioLibrary}.

        \item \texttt{ioFtdiError} \\
        Responsibility: Defines FTDI-specific library failures.

        \item \texttt{ioStreamLifecycle} \\
        Responsibility: Defines open and close behavior for stream implementations.

        \item \texttt{ioReadableByteStream} \\
        Responsibility: Defines the read-side byte-stream protocol.

        \item \texttt{ioWritableByteStream} \\
        Responsibility: Defines the write-side byte-stream protocol.

        \item \texttt{ioAbstractReadableByteStream} \\
        Responsibility: Shares readable-stream lifecycle and validation behavior.

        \item \texttt{ioAbstractWritableByteStream} \\
        Responsibility: Shares writable-stream lifecycle and validation behavior.

        \item \texttt{ioAbstractSessionBackedByteStream} \\
        Responsibility: Shares FTDI-session-backed stream behavior.

        \item \texttt{ioAbstractFileBackedByteStream} \\
        Responsibility: Shares file-backed stream behavior.

        \item \texttt{ioFtdiByteStream} \\
        Responsibility: Reads bytes from an FTDI device.

        \item \texttt{ioFtdiOutputByteStream} \\
        Responsibility: Writes bytes to an FTDI device.

        \item \texttt{ioFileInputByteStream} \\
        Responsibility: Reads bytes from a file.

        \item \texttt{ioFileByteStream} \\
        Responsibility: Writes bytes to a file.

        \item \texttt{ioFtdiSession} \\
        Responsibility: Owns the D2XX session lifecycle and FTDI configuration boundary.
    \end{itemize}

    \subsection{Multithreaded Transfer Pipeline Classes}
    \begin{itemize}
        \item \texttt{ioPipelineConfig} \\
        Responsibility: Stores end-to-end transfer settings for the read and write pipeline.

        \item \texttt{ioAcquisitionConfig} \\
        Responsibility: Stores read-worker timing and chunk-size settings.

        \item \texttt{ioTransferConfig} \\
        Responsibility: Stores write-worker timing and chunk-size settings.

        \item \texttt{ioPipelineController} \\
        Responsibility: Builds and coordinates the read worker, circular buffer, and write worker.

        \item \texttt{ioDataBuffer} \\
        Responsibility: Implements the required bounded thread-safe circular buffer.

        \item \texttt{ioByteCountMonitorBase} \\
        Responsibility: Shares byte-count and throughput timing behavior.

        \item \texttt{ioThreadedWorkerBase} \\
        Responsibility: Shares worker thread lifecycle behavior.

        \item \texttt{ioAbstractThreadedStreamWorker} \\
        Responsibility: Shares threaded stream-worker behavior for the read and write controllers.

        \item \texttt{ioUsbReadController} \\
        Responsibility: Reads bytes from a stream and pushes them into the circular buffer.

        \item \texttt{ioUsbWriteController} \\
        Responsibility: Pops bytes from the circular buffer and writes them to a stream.

        \item \texttt{ioRecoveryManager} \\
        Responsibility: Coordinates safe-stop recovery behavior across worker threads.

        \item \texttt{ioAcquisitionMonitor} \\
        Responsibility: Tracks read counts and read throughput.

        \item \texttt{ioThroughputMonitor} \\
        Responsibility: Tracks write counts and write throughput.

        \item \texttt{ioRateSchedulerBase} \\
        Responsibility: Shares fixed-rate scheduling behavior.

        \item \texttt{ioInputScheduler} \\
        Responsibility: Applies read-side timing policy.

        \item \texttt{ioOutputScheduler} \\
        Responsibility: Applies write-side timing policy.
    \end{itemize}

    \subsection{Legacy Compatibility Classes}
    \begin{itemize}
        \item \texttt{ioBuffer} \\
        Responsibility: Preserves the older fixed-size buffer abstraction used by the original FTDI operation path.

        \item \texttt{ioBaseIoOperation} \\
        Responsibility: Shares the older timed FTDI operation behavior.

        \item \texttt{ioRead} \\
        Responsibility: Implements the legacy FTDI read operation over \texttt{ioBuffer}.

        \item \texttt{ioWrite} \\
        Responsibility: Implements the legacy FTDI write operation over \texttt{ioBuffer}.
    \end{itemize}

    \section{Namespace and File Organization \\
    \small{\textit{-- Joris Wilson, Eva Kemp, and Nica Saoi}}}
    \index{Section!Namespace and File Organization}
    \label{Section::Namespace and File Organization}

    \subsection{Namespace Convention}
    The canonical implementation follows the required \texttt{io<ClassName>} naming pattern. Examples include \texttt{ioOscilloscopeController}, \texttt{ioOscilloscopeModel}, \texttt{ioOscilloscopeView}, \texttt{ioDataBuffer}, and \texttt{ioFtdiDevice}.

    Compatibility aliases such as \texttt{OscilloscopeController}, \texttt{DataBuffer}, and \texttt{FtdiDevice} are still present for older demos and tests, but they are not the primary submission path. The current architecture, diagrams, and professor-facing documentation use the canonical \texttt{io...} class names.

    \subsection{One Class Per File}
    The canonical implementation follows the one-class-per-file rule for the actual class definitions used in the architecture. The repository also contains re-export modules and compatibility surfaces, but those files do not introduce additional primary classes and should be interpreted as compatibility infrastructure rather than violations of the design rule.

    \section{Two Selectable User Interfaces \\
    \small{\textit{-- Joris Wilson, Eva Kemp, and Nica Saoi}}}
    \index{Section!Two Selectable User Interfaces}
    \label{Section::Two Selectable User Interfaces}

    \subsection{View 1}
    The first selectable Qt user interface is \texttt{ioCompactOscilloscopeView}. It presents a compact waveform display with a short status summary and is intended for fast interaction and configuration review.

    \subsection{View 2}
    The second selectable Qt user interface is \texttt{ioDetailedOscilloscopeView}. It presents a larger waveform area and a richer telemetry panel for more detailed inspection of the current oscilloscope session.

    \subsection{View Selection Mechanism}
    View selection is managed by \texttt{ioQtScopeWindow}. The window hosts concrete instances of both Qt views, switches the active visible view at runtime, and forwards view-selection and control changes to \texttt{ioOscilloscopeController}.

    \subsection{Maintaining One View Abstraction}
    The design intentionally uses one abstract View boundary only. Both selectable Qt user interfaces inherit from \texttt{ioOscilloscopeView}, and the host window is documented as a shell around those views rather than a second View abstraction.

    \section{Pipeline Processing for Scale and Shift \\
    \small{\textit{-- Joris Wilson, Eva Kemp, and Nica Saoi}}}
    \index{Section!Pipeline Processing for Scale and Shift}
    \label{Section::Pipeline Processing for Scale and Shift}

    \subsection{Scale Stage}
    The scale stage is implemented by \texttt{ioScaleFilter}. It multiplies each sample by the configured scale factor before rendering.

    \subsection{Shift Stage}
    The shift stage is implemented by \texttt{ioOffsetFilter}. It adds the configured offset value to each sample before rendering.

    \subsection{Pipeline Data Flow}
    Pipeline processing is organized through \texttt{ioFilterPipeline}, which owns the ordered list of filter stages. The model rebuilds or reuses the pipeline as controls change, and the processed result is then converted into an immutable \texttt{ioRenderState} for the views.

    \section{FTDI Communication and Circular Buffers \\
    \small{\textit{-- Joris Wilson, Eva Kemp, and Nica Saoi}}}
    \index{Section!FTDI Communication and Circular Buffers}
    \label{Section::FTDI Communication and Circular Buffers}

    \subsection{FTDI Access}
    FTDI access is handled through a dedicated backend layer. The session layer is implemented by \texttt{ioFtdiSession}, and the stream-facing FTDI endpoints are implemented by \texttt{ioFtdiByteStream} and \texttt{ioFtdiOutputByteStream}.

    The oscilloscope package itself consumes FTDI input through \texttt{ioFtdiWaveformSource} and live-session helpers rather than importing driver details into the View layer. This keeps FTDI-specific behavior localized and testable.

    \subsection{Read and Write Threads}
    The multithreaded transfer structure uses separate worker threads for reading and writing. The read thread is implemented by \texttt{ioUsbReadController}, and the write thread is implemented by \texttt{ioUsbWriteController}.

    These two worker classes are coordinated by \texttt{ioPipelineController}. Their loops are regulated by \texttt{ioInputScheduler} and \texttt{ioOutputScheduler}, while status and failure handling are delegated to the monitor and recovery classes.

    \subsection{Circular Buffer Integration}
    The circular buffer integration is centered on \texttt{ioDataBuffer}. The read worker pushes newly acquired bytes into the shared circular buffer, and the write worker removes buffered bytes from that same buffer for file or FTDI output.

    The current branch implements the required one-buffer producer-consumer design. The extra-credit two-device and two-buffer extension is documented as future work rather than as a completed feature in this branch.

    \subsection{Live Acquisition Integration}
    In live oscilloscope mode, \texttt{ioLiveOscilloscopeSession} creates the input stream, the optional tee-output stream, the transfer configuration, and the \texttt{ioPipelineController}. This keeps the live sampling path aligned with the assignment's multithreaded circular-buffer requirement while still preserving MVC separation.
\providecommand{\ioArchitectureSvgFigure}[4]{
\begin{figure}[H]
    \centering
    \IfFileExists{svg/#1}{
        \includesvg[width=#4\linewidth]{svg/#1}
    }{
        \fbox{
            \parbox{0.82\textwidth}{
                \centering
                \textbf{Diagram asset missing.} Expected SVG file: \texttt{svg/#1}.
            }
        }
    }
    \caption{#2}
    \label{#3}
\end{figure}
}

\section{4+1 Diagram Set \\
\small{\textit{-- Joris Wilson, Eva Kemp, and Nica Saoi}}}
\index{Section!4+1 Diagram Set}
\label{Section::4+1 Diagram Set}

The architecture diagram set is documented with the submitted SVG assets only. Figures~\ref{fig:qt_mvc}, \ref{fig:signal_flow}, \ref{fig:logical_view}, \ref{fig:process_view}, \ref{fig:deployment_view}, \ref{fig:activity_view}, \ref{fig:use_case_view}, \ref{fig:scenario_view_8_1}, and \ref{fig:scenario_view_8_2} together cover the Qt MVC structure, signal flow, logical grouping, process behavior, deployment structure, activity flow, use cases, and two scenario sequences.

\subsection{Qt MVC Diagram}
\index{Section!Qt MVC Diagram}
\label{Section::Qt MVC Diagram}

\ioArchitectureSvgFigure{1_qt_mvc.svg}{Qt MVC class slice}{fig:qt_mvc}{1.0}

Figure~\ref{fig:qt_mvc} is the primary Qt MVC class slice. It explains the controller, model, view hierarchy, host window, render-state handoff, theme hierarchy, and live-session attachment used by the Qt oscilloscope path.

\begin{enumerate}
    \item This diagram is the primary Qt MVC class slice. \texttt{ioOscilloscopeController} is the controller class, \texttt{ioOscilloscopeModel} is the model class, and \texttt{ioOscilloscopeView} is the single abstract View class for the Qt path.
    \item The two selectable Qt user interfaces are modeled through inheritance. \texttt{ioCompactOscilloscopeView} and \texttt{ioDetailedOscilloscopeView} are concrete subclasses of the abstract class \texttt{ioOscilloscopeView}.
    \item \texttt{ioQtScopeWindow} is a host-shell class, not a second View abstraction. It owns view-instance selection and visibility, then forwards user actions into the controller.
    \item The model class aggregates state and processing context. \texttt{ioControlState} and \texttt{ioViewportState} hold user-facing state, while \texttt{ioFilterPipeline} and \texttt{ioWaveformGenerator} stay behind the model boundary.
    \item The render handoff is made explicit through \texttt{ioRenderState}. That keeps rendering data separate from controller orchestration and helps the View layer stay presentation-focused.
    \item The theme path is also an inheritance hierarchy. \texttt{ioViewTheme} is the abstract contract, \texttt{ioStaticThemeBase} is an abstract shared base, and \texttt{ioPortraitTheme} plus \texttt{ioLandscapeTheme} are concrete subclasses.
    \item \texttt{ioLiveOscilloscopeSession} is attached to the controller rather than the View. That keeps live acquisition as application behavior instead of UI behavior.
\end{enumerate}

\subsection{Signal Flow Diagram}
\index{Section!Signal Flow Diagram}
\label{Section::Signal Flow Diagram}

\ioArchitectureSvgFigure{2_signal_flow.svg}{Signal flow and live acquisition class slice}{fig:signal_flow}{1.0}

Figure~\ref{fig:signal_flow} narrows the architecture to live acquisition, waveform source selection, byte-stream abstraction, and filter-pipeline processing.

\begin{enumerate}
    \item This class slice shows how live acquisition and waveform processing stay behind the MVC boundary. \texttt{ioOscilloscopeController} owns session startup and shutdown, while \texttt{ioOscilloscopeModel} owns waveform data and the active processing pipeline.
    \item \texttt{ioLiveOscilloscopeSession} is the coordinating class for live acquisition. It owns the pipeline-related helper classes and keeps that behavior outside the View layer.
    \item The filter design is modeled with both interface and inheritance. \texttt{ioSignalFilter} is the interface, \texttt{ioSampleMappingFilterBase} is the abstract class for shared mapping logic, and \texttt{ioScaleFilter} plus \texttt{ioOffsetFilter} are concrete subclasses.
    \item \texttt{ioFilterPipeline} depends on the filter interface rather than on one fixed concrete filter class. That means the processing chain is extensible without rewriting the pipeline class.
    \item Source selection follows the same UML pattern. \texttt{ioSignalSource} is the interface, and \texttt{ioGeneratedWaveformSource}, \texttt{ioFileWaveformSource}, and \texttt{ioFtdiWaveformSource} are alternative concrete implementations.
    \item \texttt{ioWaveformGenerator} acts as the selector for the active signal-source instance. It chooses which source participates at runtime while the model stays at the abstraction level.
    \item The byte-stream side is expressed through interfaces as well. \texttt{ioReadableByteStream} and \texttt{ioWritableByteStream} define the stream contracts, while helper classes such as \texttt{ioTappedReadableByteStream} and \texttt{ioNullWritableByteStream} participate around those abstractions.
    \item \texttt{ioFtdiWaveformSource} is the concrete class that bridges signal acquisition into the FTDI stream layer by depending on \texttt{ioFtdiByteStream}.
\end{enumerate}

\subsection{Logical View Diagram}
\index{Section!Logical View Diagram}
\label{Section::Logical View Diagram}

\ioArchitectureSvgFigure{3_logical_view.svg}{Logical architecture overview}{fig:logical_view}{1.0}

Figure~\ref{fig:logical_view} is the top-level logical overview. It groups the system into architectural subsystems and shows the dependency direction between those groups.

\begin{enumerate}
    \item This is the top-level logical overview, so it is intentionally more architectural than class-by-class. The rounded rectangles are subsystem groupings rather than full detailed class declarations.
    \item The bootstrap group contains the entry-point classes and verification helpers, including \texttt{ioScopeShell}, \texttt{ioPipelineCli}, \texttt{ioLegacyFtdiCli}, and \texttt{ioFileComparator}.
    \item The Qt MVC group is the primary user-interface subsystem. Its detailed abstract-class and inheritance structure is intentionally moved into the companion MVC slice so this overview stays readable.
    \item The signal-flow group and live-acquisition group are separated on purpose. One group handles waveform-processing strategy, while the other group handles live-session behavior and tee-history capture.
    \item The transfer-pipeline group represents the required producer-consumer design. It connects the live acquisition subsystem to the FTDI stream subsystem through the multithreaded buffer path.
    \item The FTDI streams and session group is the hardware-facing layer. It contains the deployed stream and device classes that the upper logical subsystems depend on.
    \item The dependency arrows communicate architectural direction between subsystems. They show who depends on whom, while the detailed class inheritance and interface realization relationships remain in the focused companion diagrams.
\end{enumerate}

\subsection{Process View Diagram}
\index{Section!Process View Diagram}
\label{Section::Process View Diagram}

\ioArchitectureSvgFigure{4_process_view.svg}{Process view}{fig:process_view}{1.0}

Figure~\ref{fig:process_view} shows the concurrent runtime structure of the oscilloscope system. It separates the Qt UI thread from the background transfer workers.

\begin{enumerate}
    \item This is a UML process view, so the emphasis shifts from classes to concurrent runtime components.
    \item The \texttt{Qt UI Thread} rectangle contains the UI-facing component instances: the host window, the two concrete view components, the controller, the model, the live session, the pipeline controller, the filter pipeline, and the waveform generator.
    \item The two background worker contexts are modeled explicitly as separate thread regions. \texttt{ioUsbReadController} lives in the read thread, and \texttt{ioUsbWriteController} lives in the write thread.
    \item \texttt{ioDataBuffer} is drawn as a queue instance because it is the shared runtime handoff point between the producer and consumer threads. This is the concrete concurrent buffer instance that makes the pipeline work.
    \item The controller component still acts as the orchestration center at runtime. It updates model state, starts and stops the live session, and dispatches rendering to the selected concrete Qt view component.
    \item The read-side and write-side stream endpoints are modeled as runtime components behind the worker controllers. That keeps the thread roles clear: one side produces bytes into the queue, and the other consumes bytes out of it.
\end{enumerate}

\subsection{Deployment View Diagram}
\index{Section!Deployment View Diagram}
\label{Section::Deployment View Diagram}

\ioArchitectureSvgFigure{5_deployment View.svg}{Deployment view}{fig:deployment_view}{1.0}

Figure~\ref{fig:deployment_view} describes how the software artifacts, runtime libraries, configuration files, display output, USB bus, and FTDI devices are deployed.

\begin{enumerate}
    \item This is the UML deployment view. The main concept here is not inheritance but deployment across nodes, artifacts, and external devices.
    \item The \texttt{Windows Host} node contains the deployed software artifacts for the project: the Python runtime, the controller artifact, the oscilloscope package, the ioLibrary package, the Qt libraries, and the FTDI driver interface.
    \item \texttt{io\_scope\_settings.json} is modeled as a database-style artifact because it persists runtime configuration, while the sample input is modeled as a file artifact because it represents an input resource rather than an executing component.
    \item The \texttt{Display} node hosts the \texttt{Qt window} artifact. That makes the UI deployment boundary explicit and separates the visible interface from the host software stack.
    \item The \texttt{USB Bus} node captures the hardware-facing side of the deployment. The FTDI read device is the primary external device artifact, and the FTDI write device is shown as an optional artifact for single-buffer transfer mode.
    \item The deployment links show which host artifacts depend on which libraries and external endpoints. For example, the oscilloscope package depends on \texttt{PySide6}, \texttt{pyqtgraph}, settings storage, file input, and the visible UI artifact.
    \item The ioLibrary package is the artifact that bridges software into the FTDI hardware layer. It depends on \texttt{ftd2xx.dll / D2XX} and communicates with the FTDI device artifacts on the USB bus.
\end{enumerate}

\subsection{Activity View Diagram}
\index{Section!Activity View Diagram}
\label{Section::Activity View Diagram}

\ioArchitectureSvgFigure{6_activity_view.svg}{Required activity view}{fig:activity_view}{1.0}

Figure~\ref{fig:activity_view} documents the required pipeline activity flow. It focuses on initialization, read-side behavior, shared buffer coordination, write-side behavior, and shutdown.

\begin{enumerate}
    \item This activity view shows the end-to-end producer-consumer behavior required by the oscilloscope transfer path.
    \item The swimlanes separate the responsibilities of \texttt{ioPipelineController}, \texttt{ioUsbReadController}, \texttt{ioDataBuffer}, and \texttt{ioUsbWriteController}.
    \item Initialization starts by parsing configuration, creating \texttt{ioDataBuffer}, selecting the readable and writable stream implementations, and starting the read and write controller threads.
    \item The reader thread chooses between file input and FTDI input, reads byte chunks from the selected source, and reacts to stop requests, exhausted input, missing bytes, and failure cases.
    \item The buffer lane shows the synchronization boundary. The reader waits when the buffer is full, pushes byte chunks when capacity is available, and closes the buffer when the writer must drain or stop.
    \item The writer thread waits when the buffer is empty, pops byte chunks when data is available, and writes those chunks to the selected output stream.
    \item Shutdown is explicit: the controller requests the reader stop, closes the buffer, requests the writer stop, joins both worker threads, and lets the writer close the selected output stream during stop.
\end{enumerate}

\subsection{Use Case Diagram}
\index{Section!Use Case Diagram}
\label{Section::Use Case Diagram}

\ioArchitectureSvgFigure{7_use_case.svg}{Required use case view}{fig:use_case_view}{1.0}

Figure~\ref{fig:use_case_view} shows the user-visible behavior of the system from the presenter actor's point of view.

\begin{enumerate}
    \item The actor is the presenter, and the system boundary is the \texttt{ioScope Demo}.
    \item The presenter can launch the Qt oscilloscope and can also launch the CLI shell. These are separate entry points because the graphical demo and command-line pipeline demo are separate user workflows.
    \item Launching the Qt oscilloscope includes switching between the two Qt user interfaces, displaying a waveform, applying the scale-and-offset pipeline, and acquiring a waveform from FTDI.
    \item FTDI acquisition connects to the multithreaded FTDI pipeline use case. That relationship keeps the visible oscilloscope behavior tied to the required backend transfer behavior.
    \item Switching between two Qt user interfaces is modeled as a user-facing capability, not as an implementation detail. This matches the design shown in Figure~\ref{fig:qt_mvc}, where the two concrete views share one abstract View contract.
    \item The use-case view intentionally avoids internal class details. Those details are covered in the class and runtime views shown in Figures~\ref{fig:qt_mvc}, \ref{fig:signal_flow}, and \ref{fig:process_view}.
\end{enumerate}

\subsection{Scenario View 8.1 Diagram}
\index{Section!Scenario View 8.1 Diagram}
\label{Section::Scenario View 8.1 Diagram}

\ioArchitectureSvgFigure{8.1_scenario view.svg}{Scenario view 8.1: start FTDI session with tee-to-file}{fig:scenario_view_8_1}{1.0}

Figure~\ref{fig:scenario_view_8_1} shows the start sequence for an FTDI live session that also tees captured bytes to a file.

\begin{enumerate}
    \item The scenario begins when the user clicks Start in the Qt window. \texttt{ioQtScopeWindow} forwards that request to \texttt{ioOscilloscopeController}.
    \item The controller refreshes samples and enters its live-session startup path before delegating to \texttt{ioLiveOscilloscopeSession}.
    \item \texttt{ioLiveOscilloscopeSession} builds the input stream from the FTDI settings, creates the \texttt{ioFtdiByteStream}, and wraps that input with \texttt{ioTappedReadableByteStream} so sample history can be captured while bytes continue through the transfer path.
    \item The same session builds the output stream from the tee settings. In this scenario, the tee output is a file stream for \texttt{demo\_ftdi\_capture.bin}.
    \item The session creates \texttt{ioPipelineConfig}, creates \texttt{ioPipelineController} with the tapped input and file output, and starts the pipeline.
    \item Starting the pipeline starts the writer thread and then the reader thread. This matches the producer-consumer activity shown in Figure~\ref{fig:activity_view}.
    \item After startup, the controller appends a session-start event, renders an \texttt{ioRenderState}, and the detailed Qt view shows FTDI plus tee-to-file status.
\end{enumerate}

\subsection{Scenario View 8.2 Diagram}
\index{Section!Scenario View 8.2 Diagram}
\label{Section::Scenario View 8.2 Diagram}

\ioArchitectureSvgFigure{8.2_scenario view.svg}{Scenario view 8.2: Qt startup, configuration, processing, and render refresh}{fig:scenario_view_8_2}{1.0}

Figure~\ref{fig:scenario_view_8_2} shows the normal Qt startup and rendering scenario, including view creation, control updates, waveform generation, filtering, and render dispatch.

\begin{enumerate}
    \item The scenario begins with \texttt{python controller.py scope-qt}. That command creates \texttt{ioQtScopeWindow} with the two concrete view instances.
    \item The controller is created with the same views, then the compact view is selected as the active view through \texttt{setActiveView("compact")}.
    \item The user chooses the source, scale, and offset. The window forwards these selections to the controller through \texttt{setInputSource(...)}, \texttt{setScale(...)}, and \texttt{setOffset(...)}.
    \item The controller updates \texttt{ioOscilloscopeModel} so that the user-facing control state stays behind the model boundary.
    \item When acquisition starts, the controller asks the model to refresh samples. The model delegates source selection to \texttt{ioWaveformGenerator}.
    \item The waveform generator returns the raw signal, and \texttt{ioFilterPipeline} processes that signal through the active scale and offset filters.
    \item The controller renders the resulting \texttt{ioRenderState} to both concrete view instances. The host window then shows the active concrete view instance.
    \item The scenario ends with both the compact and detailed UI instances updated. This reinforces the design in Figure~\ref{fig:qt_mvc}: there are two concrete Qt interfaces, but they are coordinated through one controller and one shared View abstraction.
\end{enumerate}

\section{Class Diagram \\
\small{\textit{-- Joris Wilson, Eva Kemp, and Nica Saoi}}}
\index{Section!Class Diagram}
\label{Section::Class Diagram}

The class architecture for this assignment is documented by the focused SVG figures in Section~\ref{Section::4+1 Diagram Set}. Figure~\ref{fig:qt_mvc} is the primary Qt MVC class slice, Figure~\ref{fig:signal_flow} covers signal processing and live acquisition classes, and Figure~\ref{fig:logical_view} provides the top-level logical grouping used to relate those slices to the rest of the system.

Together, these class-oriented diagrams show the required one abstract Qt View boundary, two concrete selectable Qt user interfaces, the model-owned processing pipeline, the controller-owned live session, and the FTDI stream abstractions used by the live acquisition path.

\section{AI Prompt and Generated Output \\
    \small{\textit{-- Joris Wilson, Eva Kemp, and Nica Saoi}}}
    \index{Section!AI Prompt and Generated Output}
    \label{Section::AI Prompt and Generated Output}

    \subsection{Prompt Given to the AI Tool}
    The AI tool was given a prompt requesting a Qt-based MVC redesign with one abstract View and two selectable concrete Qt user interfaces. The prompt required one-class-per-file organization, canonical \texttt{io<ClassName>} naming, preservation of the existing scale-and-shift pipe-and-filter design, preservation of the multithreaded circular-buffer FTDI transfer path, and use of the earlier \texttt{origin/oscilloscope} branch as historical source material rather than as the canonical final architecture.

    \subsection{Generated Output Summary}
    The generated output established the current Qt architecture direction by introducing the abstract View boundary, the two concrete Qt user interfaces, the Qt host window, and the supporting controller and model updates needed to preserve MVC behavior. Additional refinement was then performed to align naming, strengthen abstraction boundaries, preserve compatibility with the existing codebase, and integrate live FTDI or file-backed acquisition through the existing backend pipeline.

    \section{Implementation Notes and Changes \\
    \small{\textit{-- Joris Wilson, Eva Kemp, and Nica Saoi}}}
    \index{Section!Implementation Notes and Changes}
    \label{Section::Implementation Notes and Changes}

    After the initial AI-generated implementation was produced, several revisions were made to align the final design with the assignment requirements and the evolving codebase. The Qt classes were made import-safe for environments where \texttt{PySide6} and \texttt{pyqtgraph} are unavailable, the controller was updated to fan out render state to both concrete views while tracking the active visible view, the live acquisition path was moved into \texttt{ioLiveOscilloscopeSession}, and the FTDI adapter surface was aligned with the \texttt{io<ClassName>} naming rule through canonical \texttt{ioFtdiDevice} naming plus preserved compatibility aliases.

    \section{Proof of Functionality \\
    \small{\textit{-- Joris Wilson, Eva Kemp, and Nica Saoi}}}
    \index{Section!Proof of Functionality}
    \label{Section::Proof of Functionality}

    The execution of the program, including the file-to-file demonstration path and the walkthrough of how the controller code coordinates with the oscilloscope, is shown in the accompanying YouTube demo submitted with the report:

    \begin{quote}
    \url{https://youtu.be/hW8-CKr5dnA}
    \end{quote}

    \begin{itemize}
        \item Switching between the compact and detailed Qt user interfaces.
        \item Displaying waveform data in the active Qt user interface.
        \item Applying scale and offset processing through the filter pipeline.
        \item Reading from FTDI or file-backed sources through the live-session path.
        \item Writing tee output to a file or FTDI sink through the multithreaded transfer pipeline.
    \end{itemize}

    \section{Conclusion \\
    \small{\textit{-- Joris Wilson, Eva Kemp, and Nica Saoi}}}
    \index{Section!Conclusion}
    \label{Section::Conclusion}

    The current oscilloscope architecture satisfies the assignment's main design goals while matching the implemented codebase and diagrams. The system preserves one abstract View boundary, provides two selectable Qt user interfaces, keeps scale and shift in the filter pipeline, isolates FTDI communication in the backend stream hierarchy, and uses a real multithreaded circular-buffer transfer path for acquisition and output.

